\section{Práctica de laboratorio}

\subsection{Identificación del sensor de oximetría}
En la primera etapa del trabajo práctico se estudió el sensor de SpO$_2$ BCI, utilizado para la medición de la saturación de oxígeno mediante oximetría de pulso. Para reconocer las conexiones eléctricas del sensor, se analizaron el cable paciente y su conector DB9 mediante mediciones de resistencia y continuidad. Estas mediciones permitieron identificar las conexiones correspondientes a los componentes del sensor para utilizarlas en el armado del circuito. 

\subsection{Armado de circuito sensor de luz roja e IR}

Durante la práctica de laboratorio se implementó en un protoboard un circuito que sensa la intensidad de la luz de los emisores rojo e IR. Su funcionamiento se basa en el funcionamiento del amplificador operacional TL082, que convierte la corriente generada por el fotodetector en una tensión de salida, adecuada para su medición en el osciloscopio. 

Los componentes utilizados para la implementación fueron:

\begin{itemize}
    \item Integrado TL082, utilizando dos de sus amplificadores operacionales.
    \item Resistencia de $1M~\Omega$ para la realimentación del TL082.
    \item Resistencia de $220~\Omega$ para modificar la intensidad de la luz del sensor.
    \item Sensor SpO$_2$ BCI con ficha DB9
    \item Protoboard y cables de conexión.
    \item 2 baterías de $9~\mathrm{V}$
\end{itemize}

En la Figura~ se muestra el circuito amplificador de transimpedancia analizado. El circuito incorpora una resistencia de realimentación de $R_2=1\,\mathrm{M}\Omega$, conectada entre la salida y la entrada inversora del amplificador que determina la relación entre la corriente de entrada y la tensión de salida.
Asimismo, se representan los diodos $D_1$ y $D_2$ presentes en el sensor de SpO$_2$ BCI, junto con una resistencia de $220\,\Omega$ asociada al diodo emisor. El amplificador operacional se alimenta con tensiones de $+9\,\mathrm{V}$ y $-9\,\mathrm{V}$, mientras que su entrada no inversora se conecta a tierra.

\begin{figure}[htbp]
    \centering
    \begin{circuitikz}[scale=1, transform shape]
    
        % Amplificador operacional
        \draw (4,0) node[op amp] (U1) {};
        \node at (4,-0.8) {TL082};
        
        % Entrada no inversora a tierra
        \draw (U1.+) -- ++(-1,0) node[ground]{};
        
        % Nodo inversor
        \coordinate (N) at (2.5,0.3);
        \draw (U1.-) -- (N);
        
        % Realimentación: R2 y C1 en paralelo
        \draw (N) -- (2.5,2.5) -- (3,2.5)
            to[R,l={$R_2=1\,\mathrm{M}\Omega$}] (5,2.5)
            -- (5,0.3);
        
        % Salida
        \draw (U1.out) -- ++(1.2,0) node[right] {$V_{\mathrm{out}}$};
        
        % Diodo D1 en la entrada inversora
        \draw (N) -- (1.2,0.3)
            to[D,l_=$D_1$] (1.2,2.0)
            node[above] {$+9\,\mathrm{V}$};
        
        % Diodo D2 con resistencia de 220 ohm en serie
        \draw (-1,1.2) node[above] {$+9\,\mathrm{V}$}
            to[D,l={$D_2$}] (-1,-0.3)
            to[R,l={$220\,\Omega$}] (-1,-1.8)
            node[ground] {};
        
        % Alimentación simétrica del operacional
        \draw (U1.up) -- ++(0,1.0) node[above] {$+9\,\mathrm{V}$};
        \draw (U1.down) -- ++(0,-1.0) node[below] {$-9\,\mathrm{V}$};
        
    \end{circuitikz}
    \caption{Circuito amplificador de transimpedancia.}
    \label{fig:transimpedancia}
\end{figure}



